\section{表面活性物质}


\begin{frame}
    \frametitle{表面活性物质的分类}

\begin{table}[]
    \begin{tabular}{cc|l}
    \Xhline{1pt}
    \rowcolor{maincolor!10}
    \multicolumn{2}{c|}{类别} & \multicolumn{1}{c}{实例} \\ \hline
    \multicolumn{1}{l|}{\multirow{6}{*}{\makecell{离子\\型\\表面\\活性\\物质}}}  &  \multirow{2}{*}{阴离子型} & 羧酸盐\ch{RCOO^-M+}，硫酸酯盐\ch{ROSO3^-M+}   \\ 
    \multicolumn{1}{l|}{}  &   &  磺酸盐\ch{RSO3^-M+}，磷酸酯盐\ch{ROPO3^-M+}                \\ \cline{2-3} 
    \multicolumn{1}{l|}{}  &  \multirow{2}{*}{阳离子型} &  伯胺盐\ch{RNH3^+X^-}，季胺盐\ch{RN^+(CH3)3X^-}\\ 
    \multicolumn{1}{l|}{}  &   &  吡啶盐                \\ \cline{2-3} 
    \multicolumn{1}{l|}{}  &  \multirow{2}{*}{两性离子型} &  氨基酸型\ch{RN^+CH2CH2COO^-}              \\ 
    \multicolumn{1}{l|}{}  &   &  甜菜碱型\ch{RN^+(CH3)2CH2COO^-}                 \\ \hline
    \multicolumn{2}{c|}{\multirow{3}{*}{非离子型表面活性物质}}  &   聚氧乙烯醚\ch{RO(CH2CH2O)_n H}   \\ \cline{3-3} 
    &   &  聚氧乙烯酯\ch{RCOO(CH2CH2O)_nH}    \\ \cline{3-3} 
    &   &  多元醇型\ch{RCOOCH2C(CH2OH)3}      \\ \Xhline{1pt}
    \end{tabular}
    \end{table}

\end{frame}


\begin{frame}
    \frametitle{表面活性物质的HLB值}

    格里芬(Griffin)提出用HLB(hydrophile-lipophile balance)，即亲水$-$亲油平衡值来表示表面活性物质的亲水性和亲油性的相对强弱。
    \begin{itemize}
      \item HLB越大，其亲水性越强；
      \item HLB越小，其亲水性越弱，亲油性越强。
      \item 表面活性物质的HLB是个相对值。
    \end{itemize}
    \begin{table}[]
      \caption{HLB范围及其适当用途}
      \begin{tabular}{cc|cc}
        \Xhline{1pt}
      \rowcolor{maincolor!10}
      HLB & 应用           & HLB   & 应用           \\ \hline
      1-3 & 消泡剂          & 8-18  & O/W(水包油)型乳化剂 \\
      3-6 & W/O(油包水)型乳化剂 & 13-15 & 洗涤剂          \\
      7-9 & 润湿剂          & 15-18 & 增溶剂          \\ \Xhline{1pt}
      \end{tabular}
      \end{table}
\end{frame}